\documentclass[11pt]{article}

\usepackage[a4paper,margin=1in]{geometry}
\usepackage{amsmath,amssymb,amsthm,mathtools}
\usepackage{graphicx}
\usepackage{microtype}
\usepackage[hidelinks]{hyperref}

\newtheorem{theorem}{Theorem}
\newtheorem{lemma}{Lemma}
\newtheorem{remark}{Remark}
\newcommand{\R}{\mathbb R}
\newcommand{\cH}{\mathcal H}

\title{An Analytic Two-Dimensional Counterexample\\
to the Wronski-Basis Conjecture}
\author{Candidate counterexample for arXiv:2508.04639\\
Run \texttt{fa\_banach\_001}, lane 7}
\date{12 August 2026}

\begin{document}
\maketitle

\begin{center}
\fbox{\parbox{0.91\textwidth}{\textbf{Status.}
Candidate full counterexample, likely valid, pending specialist human review.
It disproves Conjecture~8 already for a two-dimensional Hilbert space of
entire functions.}}
\end{center}

\section{Source question}

Micheas asks whether every orthonormal basis of a Hilbert function space is
the fundamental solution set, via Wronski orthogonalization, of a unique
linear homogeneous ordinary differential equation (HODE), possibly of
infinite order \cite[Conjecture~8, p.~16]{Micheas}.

\begin{figure}[ht]
  \centering
  \includegraphics[width=\textwidth]{figures/open_problem_crop.png}
  \caption{Conjecture~8 in arXiv:2508.04639v4, PDF page 16.}
\end{figure}

Immediately after the conjecture, the paper says that it is satisfied
``trivially'' for a finite basis $\{e_1,\ldots,e_N\}$ and identifies the
claimed equation as the $N$th-order Wronskian quotient
\[
  \frac{W(e_1,\ldots,e_N,y)}{W(e_1,\ldots,e_N)}=0.
\]
The paper's regular-ODE convention requires continuous coefficients on the
domain and a leading coefficient that is nonzero at every point
\cite[Remark~1]{Micheas}.  We use exactly that finite-dimensional
interpretation.

\section{Proof intuition}

Linear independence of functions does not imply that their Wronskian is
pointwise nonzero.  This distinction is invisible if one divides formally by
the Wronskian.  The smallest possible example is $1,x^2$: they are linearly
independent and can be made an orthonormal basis, but
$W(1,x^2)=2x$ vanishes at the origin.  The formal Wronskian equation is
$xy''-y'=0$, which is necessarily singular at that point.  Direct
substitution shows that no other regular second-order equation can avoid the
same obstruction.

\section{Counterexample}

Let
\[
  \cH=\operatorname{span}_{\R}\{1,x^2\}
\]
as a space of real-valued functions on $\R$, and define
\[
 \big\langle a+bx^2,c+dx^2\big\rangle_{\cH}=ac+bd.
\]

\begin{theorem}\label{thm:counterexample}
The space $\cH$ is a Hilbert function space and
$\mathcal E=\{1,x^2\}$ is an orthonormal basis of entire functions.  There
is no regular linear second-order HODE on $\R$ having $\mathcal E$ as a
fundamental set of solutions.  Consequently, Conjecture~8 and its asserted
finite-basis case are false as stated.
\end{theorem}

\begin{proof}
The coefficient map
\[
  \R^2\longrightarrow\cH,\qquad (a,b)\longmapsto a+bx^2,
\]
is an isometric linear bijection.  Thus $\cH$ is complete, and
$\{1,x^2\}$ is an orthonormal basis.  Both functions are entire.

Suppose that a regular second-order linear HODE on $\R$,
\begin{equation}\label{eq:general}
 p_2(x)y''(x)+p_1(x)y'(x)+p_0(x)y(x)=0,
\end{equation}
has $1$ and $x^2$ as solutions.  Here the coefficients are continuous and
$p_2(x)\ne0$ for every $x\in\R$, as required by the source paper's regularity
convention.  Substituting $y=1$ into \eqref{eq:general} gives
\[
  p_0(x)=0\qquad(x\in\R).
\]
Substituting $y=x^2$ then gives
\[
  2p_2(x)+2x p_1(x)=0\qquad(x\in\R).
\]
At $x=0$ this says $p_2(0)=0$, contradicting regularity.  Hence no such HODE
exists.

A fundamental set for a regular $N$th-order equation has $N$ elements; this
is also the convention used in the source paper's finite-basis paragraph.
Thus the two-element basis is precisely the $N=2$ case claimed there, and
failure of existence already disproves the universal conjecture.
\end{proof}

\section{The failed Wronskian formula}

The source paper's proposed formula makes the singularity completely
explicit:
\[
 W(1,x^2)=
 \det\begin{pmatrix}1&x^2\\0&2x\end{pmatrix}=2x,
\]
whereas expansion of the augmented determinant gives
\[
 W(1,x^2,y)=
 \det\begin{pmatrix}
  1&x^2&y\\
  0&2x&y'\\
  0&2&y''
 \end{pmatrix}
 =2xy''-2y'.
\]
Away from the origin, division by $W(1,x^2)$ produces
$y''-x^{-1}y'=0$.  On all of $\R$, its undivided form
$xy''-y'=0$ has leading coefficient $x$, so it is singular exactly where the
denominator vanishes.  Reordering or rescaling the basis only multiplies the
Wronskian by a nonzero constant and cannot repair this zero.

More generally, if two $C^2$ functions form a fundamental system for a
regular second-order equation on an interval, Abel's identity forces their
Wronskian either to vanish nowhere or to vanish identically.  Therefore any
linearly independent pair with a Wronskian having an isolated zero supplies
the same kind of obstruction.

\section{Verification notes}

There are three independent exact checks.
\begin{enumerate}
  \item The coefficient map above proves completeness and orthonormality.
  \item Direct substitution into an arbitrary regular equation forces
        $p_2(0)=0$.
  \item Direct determinant expansion yields
        $W(1,x^2)=2x$ and $W(1,x^2,y)=2xy''-2y'$.
\end{enumerate}
No numerical computation or unproved external lemma is needed.

\section{Scope, repair, and novelty check}

The counterexample concerns the conjecture exactly as stated on the global
domain $\R$ and its own finite-order interpretation.  It does not rule out a
corrected theorem on an interval avoiding the origin: on $(0,\infty)$, for
example, $y''-x^{-1}y'=0$ is regular.  Nor does it address a newly defined
notion of ``fundamental set'' for infinite-order equations.  A necessary
finite-dimensional repair is to require the basis Wronskian to be nonzero
everywhere on the domain; under the usual differentiability assumptions,
the standard Wronskian formula then applies.

On 12 August 2026, bounded arXiv searches for the exact title, the exact
phrase ``Bases as solutions to HODEs'', and ``Wronski basis'' found the source
paper but no later arXiv paper or note answering Conjecture~8.  This is a
bounded novelty check, not a claim of exhaustive bibliographic priority.

\section{Human review recommendation}

Review as a likely valid full counterexample.  The two points worth checking
are that finite-dimensional function spaces are included in the source's
phrase ``any Hilbert function space'' (as its finite-basis paragraph itself
indicates), and that its stated regularity convention is imposed globally on
$\R$.  Once those are accepted, the contradiction is elementary.

\begin{thebibliography}{9}
\bibitem{Micheas}
Athanasios Christou Micheas,
\emph{Construction of linearly independent and orthogonal functions in
Hilbert function spaces via Wronski determinants},
arXiv:2508.04639v4, 2026.
\url{https://arxiv.org/abs/2508.04639}.
\end{thebibliography}

\end{document}
